% =============================================================================
% Aula 04 — PUT · Regra de negócio · Vendas no StockSales
% Beamer + tema Madrid | Curso Entra21 — Spring Framework
% Compilar: tectonic -X compile slides-aula-04.tex
% =============================================================================
\documentclass[aspectratio=169,11pt]{beamer}

\usetheme{Madrid}
\usecolortheme{default}

\usepackage[T1]{fontenc}
\usepackage[utf8]{inputenc}
\usepackage[brazil]{babel}
\usepackage{booktabs}
\usepackage{graphicx}
\usepackage{hyperref}
\usepackage{listings}
\usepackage{xcolor}
\usepackage{tikz}
\usetikzlibrary{arrows.meta, positioning, shapes.geometric}

\definecolor{codebg}{RGB}{245,245,245}
\definecolor{codegreen}{RGB}{40,120,40}
\definecolor{codeblue}{RGB}{30,70,160}
\definecolor{codered}{RGB}{180,50,50}

\lstset{
  basicstyle=\ttfamily\small,
  backgroundcolor=\color{codebg},
  frame=single,
  rulecolor=\color{black!20},
  breaklines=true,
  columns=fullflexible,
  keepspaces=true,
  showstringspaces=false,
  keywordstyle=\color{codeblue}\bfseries,
  commentstyle=\color{codegreen},
  stringstyle=\color{codered},
  tabsize=2,
  literate={á}{{\'a}}1 {é}{{\'e}}1 {í}{{\'i}}1 {ó}{{\'o}}1 {ú}{{\'u}}1
           {Á}{{\'A}}1 {É}{{\'E}}1 {Í}{{\'I}}1 {Ó}{{\'O}}1 {Ú}{{\'U}}1
           {ã}{{\~a}}1 {õ}{{\~o}}1 {Ã}{{\~A}}1 {Õ}{{\~O}}1
           {ç}{{\c{c}}}1 {Ç}{{\c{C}}}1
           {ê}{{\^e}}1 {Ê}{{\^E}}1 {â}{{\^a}}1 {Â}{{\^A}}1
           {à}{{\`a}}1 {À}{{\`A}}1
}

\newcommand{\onde}[1]{%
  \begin{block}{Onde estamos}
    #1
  \end{block}
}

\title{Aula 04 --- Spring Framework}
\subtitle{PUT · Regra de negócio · Vendas}
\author{Mauricio Duailibi Neto}
\institute{Turma Java Entra21}
\date{}

\begin{document}

\begin{frame}
  \titlepage
\end{frame}

\begin{frame}{Como será a aula de hoje}
  \begin{block}{Parte 1 --- 18h30 às 20h}
    \begin{itemize}
      \item Revisão da Aula 03
      \item PUT: alterar um cadastro
      \item Buscar um item pelo id
      \item Retirar quantidade, com regra
      \item Projeto \textbf{laboratorio} --- fazemos juntos
    \end{itemize}
  \end{block}

  \vspace{0.4em}
  \begin{block}{Parte 2 --- 20h20 às 22h}
    \begin{itemize}
      \item Voltamos ao \textbf{StockSales}
      \item PUT de produto, clientes e venda
    \end{itemize}
  \end{block}
\end{frame}

\begin{frame}{Dois projetos nesta aula}
  Não vamos misturar os códigos.

  \vspace{0.6em}
  \begin{description}
    \item[laboratorio] Primeira parte da aula.
      Aqui aprendemos PUT e a regra de retirar quantidade.
    \item[StockSales] O sistema da loja, que já existe.
      Só depois do intervalo.
  \end{description}

  \vspace{0.6em}
  Quando aparecer um código, o slide vai dizer:\\
  \textbf{qual projeto} e \textbf{qual arquivo} estamos mexendo.
\end{frame}

\begin{frame}{Mapa da aula}
  \tableofcontents
\end{frame}

% #############################################################################
\section{Revisão da Aula 03}

\begin{frame}{O que já sabemos}
  \begin{itemize}
    \item A lista fica \textbf{na classe} do Controller
    \item GET devolve a lista
    \item POST lê o JSON do corpo e adiciona na lista
    \item DELETE usa o id no endereço e remove o item
    \item O JavaScript chama a API com Axios
  \end{itemize}

  \vspace{0.7em}
  Faltou o PUT: \textbf{alterar} um cadastro que já existe.
\end{frame}

\begin{frame}{Os quatro verbos}
  \begin{center}
  \begin{tabular}{lll}
    \toprule
    Ação & Verbo HTTP & O que faz \\
    \midrule
    Listar / ler & GET & Lê, não muda nada \\
    Cadastrar & POST & Cria um item novo \\
    Editar & PUT & Troca os dados de um item \\
    Excluir & DELETE & Apaga um item \\
    \bottomrule
  \end{tabular}
  \end{center}

  \vspace{0.7em}
  CRUD continua o mesmo nome:\\
  Create, Read, Update, Delete.\\
  O Update é o PUT de hoje.
\end{frame}

\begin{frame}{O id na URL}
  DELETE já usou o id no endereço:

  \vspace{0.5em}
  \texttt{DELETE /api/materiais/2}

  \vspace{0.5em}
  O \texttt{2} é o \texttt{@PathVariable}.\\
  O método recebe esse número e procura na lista.

  \vspace{0.6em}
  PUT usa o mesmo jeito:\\
  \texttt{PUT /api/materiais/2}\\
  Só que, além do id, vem um \textbf{corpo} com os dados novos.
\end{frame}

% #############################################################################
\section{Projeto laboratorio}

\begin{frame}{O projeto laboratorio}
  A partir de hoje os exercícios da primeira parte\\
  ficam no projeto \textbf{laboratorio}.

  \vspace{0.5em}
  Ele já está no repositório.\\
  Não vamos gerar outro projeto no \texttt{start.spring.io}.

  \vspace{0.6em}
  \begin{block}{O que já vem pronto}
    Cadastro de materiais com GET, POST e DELETE.\\
    Lista na classe. Formulário na tela. Botão Excluir.
  \end{block}
\end{frame}

\begin{frame}{Passo 1 --- Atualizar o repositório}
  No terminal, dentro da pasta do repositório:

  \vspace{0.5em}
  \texttt{git pull}

  \vspace{0.6em}
  Deve aparecer a pasta \texttt{laboratorio/}.

  \vspace{0.5em}
  Abram essa pasta no VS Code.
\end{frame}

\begin{frame}{Passo 2 --- Subir o Spring}
  \onde{Projeto \textbf{laboratorio}}

  No terminal, dentro da pasta \texttt{laboratorio}:

  \vspace{0.4em}
  \texttt{./mvnw spring-boot:run}\\
  {\small No Windows: \texttt{.\textbackslash mvnw.cmd spring-boot:run}}

  \vspace{0.5em}
  Abram no navegador: \texttt{http://localhost:8080}

  \vspace{0.5em}
  \begin{alertblock}{Um Spring por vez}
    Se outro projeto ainda estiver na porta 8080,
    parem o outro antes.
  \end{alertblock}
\end{frame}

\begin{frame}{Quais arquivos já existem}
  Projeto: \textbf{laboratorio}

  \vspace{0.4em}
  \begin{center}
  \begin{tabular}{ll}
    \toprule
    Arquivo & Para quê \\
    \midrule
    \texttt{Material.java} & Os dados de um material \\
    \texttt{MaterialController.java} & GET, POST e DELETE \\
    \texttt{static/index.html} & A página \\
    \texttt{static/js/app.js} & O JavaScript da página \\
    \bottomrule
  \end{tabular}
  \end{center}

  \vspace{0.5em}
  Hoje vamos \textbf{acrescentar} métodos e botões.\\
  Não criamos outro projeto.
\end{frame}

\begin{frame}[fragile]{Onde cada arquivo fica}
\begin{lstlisting}[basicstyle=\ttfamily\footnotesize]
laboratorio/
  src/main/java/com/entra21/laboratorio/
    LaboratorioApplication.java
    model/Material.java
    controller/MaterialController.java
  src/main/resources/static/
    index.html
    js/app.js
\end{lstlisting}

  Java fica em \texttt{java/}.\\
  Página e JS ficam em \texttt{resources/static/}.
\end{frame}

\begin{frame}{Conferindo o que já funciona}
  \begin{itemize}
    \item [ ] A página abre em \texttt{http://localhost:8080}
    \item [ ] Caneta, Caderno e Grampeador aparecem
    \item [ ] Cadastrar um material novo atualiza a lista
    \item [ ] Excluir some da lista
    \item [ ] Recarregar a página mantém o que estava na memória
  \end{itemize}

  \vspace{0.5em}
  Se isso não estiver ok, não sigam para o PUT.\\
  A lista precisa estar na classe, como na aula passada.
\end{frame}

\begin{frame}[fragile]{Material.java --- o que a classe guarda}
  \onde{\texttt{model/Material.java}}

  Três dados: id, nome e quantidade.

\begin{lstlisting}[language=Java,basicstyle=\ttfamily\footnotesize]
public class Material {
  private int id;
  private String nome;
  private int quantidade;
}
\end{lstlisting}

  Já existem construtor vazio, construtor cheio, gets e sets.\\
  Não precisamos mexer nesta classe agora.
\end{frame}

% #############################################################################
\section{PUT}

\begin{frame}{Cadastrar e alterar não são a mesma ação}
  \begin{itemize}
    \item POST cria um material \textbf{novo} e gera um id
    \item PUT pega um material \textbf{que já existe} e troca
          nome ou quantidade
  \end{itemize}

  \vspace{0.6em}
  Se usarmos POST para editar, nasce um segundo cadastro.\\
  O antigo continua lá.

  \vspace{0.5em}
  Por isso o PUT leva o id: o servidor precisa saber
  \textbf{qual} linha da lista vai mudar.
\end{frame}

\begin{frame}[fragile]{O pedido de alteração}
  Exemplo: a Caneta passou a ter quantidade 20.

  \vspace{0.4em}
  \begin{block}{Pedido}
    \begin{itemize}
      \item Verbo: \texttt{PUT}
      \item Endereço: \texttt{/api/materiais/1}
      \item Corpo:
    \end{itemize}
\begin{lstlisting}[basicstyle=\ttfamily\footnotesize]
{"nome": "Caneta", "quantidade": 20}
\end{lstlisting}
  \end{block}

  \vspace{0.2em}
  O id vai na URL.\\
  Nome e quantidade novos vão no corpo.
\end{frame}

\begin{frame}{O que o PUT faz, em ordem}
  \begin{enumerate}
    \item Lê o id da URL
    \item Lê nome e quantidade do corpo
    \item Percorre a lista
    \item Quando o id bater: chama \texttt{setNome} e
          \texttt{setQuantidade}
    \item Devolve o material já atualizado
  \end{enumerate}

  \vspace{0.5em}
  A lista é a mesma.\\
  Só os campos daquele objeto mudam.
\end{frame}

\begin{frame}[fragile]{PUT no Controller}
  \onde{Projeto \textbf{laboratorio}\\
        Arquivo \texttt{controller/MaterialController.java}\\
        método novo, dentro da classe}

  Imports novos: \texttt{PutMapping}.\\
  \texttt{PathVariable} e \texttt{RequestBody} vocês já usam.
\end{frame}

\begin{frame}[fragile]{MaterialController.java --- atualizar}
  \onde{\texttt{MaterialController.java}}

\begin{lstlisting}[language=Java,basicstyle=\ttfamily\footnotesize]
@PutMapping("/api/materiais/{id}")
public Material atualizar(@PathVariable int id,
                          @RequestBody Material dados) {
  for (int i = 0; i < materiais.size(); i++) {
    if (materiais.get(i).getId() == id) {
      materiais.get(i).setNome(dados.getNome());
      materiais.get(i).setQuantidade(dados.getQuantidade());
      return materiais.get(i);
    }
  }
  return null;
}
\end{lstlisting}
\end{frame}

\begin{frame}{Olhando o \texttt{for} com calma}
  \begin{enumerate}
    \item \texttt{i} começa em 0 e anda até o fim da lista
    \item \texttt{materiais.get(i)} é o material daquela posição
    \item Se o id for o que veio na URL, achamos o certo
    \item \texttt{setNome} e \texttt{setQuantidade} gravam
          o que veio no JSON
    \item \texttt{return} devolve esse material e encerra o método
  \end{enumerate}

  \vspace{0.5em}
  Se o \texttt{for} terminar sem achar, devolvemos \texttt{null}.
\end{frame}

\begin{frame}{POST cria. PUT não cria}
  No POST vocês faziam:

  \vspace{0.3em}
  \texttt{material.setId(proximoId);} e \texttt{materiais.add(...)}

  \vspace{0.5em}
  No PUT \textbf{não} geramos id novo e \textbf{não} damos add.\\
  O objeto já está na lista. Só mudamos os campos.

  \vspace{0.5em}
  Se derem add no PUT, o material antigo continua
  e ainda nasce um duplicado.
\end{frame}

\begin{frame}{Editar pela tela}
  \onde{Projeto \textbf{laboratorio}\\
        Arquivo \texttt{static/js/app.js}}

  Objetivo: clicar em Editar, o formulário se preenche,
  e o envio passa a ser um PUT.
\end{frame}

\begin{frame}[fragile]{Uma variável para o id em edição}
  \onde{\texttt{static/js/app.js}, no topo do arquivo}

  Se estiver vazia, o form cadastra.\\
  Se tiver um número, o form altera.

\begin{lstlisting}[basicstyle=\ttfamily\footnotesize]
var idEditando = "";
\end{lstlisting}

  \vspace{0.4em}
  Vamos usar essa variável em dois lugares:
  \begin{itemize}
    \item no botão Editar, para guardar o id
    \item no submit do form, para escolher POST ou PUT
  \end{itemize}
\end{frame}

\begin{frame}[fragile]{Botão Editar dentro do \texttt{for}}
  \onde{\texttt{static/js/app.js}, no \texttt{for} de
        \texttt{carregarMateriais}}

\begin{lstlisting}[basicstyle=\ttfamily\scriptsize]
var btnEditar = document.createElement("button");
btnEditar.textContent = "Editar";
btnEditar.setAttribute("data-id", material.id);
btnEditar.setAttribute("data-nome", material.nome);
btnEditar.setAttribute("data-quantidade", material.quantidade);

btnEditar.addEventListener("click", function (evento) {
  idEditando = evento.target.getAttribute("data-id");
  document.getElementById("nome").value =
    evento.target.getAttribute("data-nome");
  document.getElementById("quantidade").value =
    evento.target.getAttribute("data-quantidade");
});
li.appendChild(btnEditar);
\end{lstlisting}

  O botão Editar fica \textbf{antes} do botão Excluir.
\end{frame}

\begin{frame}{Por que \texttt{data-id} e \texttt{evento.target}}
  O \texttt{for} roda rápido e termina.\\
  O clique acontece depois.

  \vspace{0.5em}
  Se o clique usar a variável \texttt{material} do \texttt{for},
  todos os botões podem pegar o último da lista.

  \vspace{0.5em}
  Por isso o id vai para o atributo do botão,
  e o clique lê \texttt{evento.target}.
\end{frame}

\begin{frame}[fragile]{O submit --- montar os dados}
  \onde{\texttt{static/js/app.js}, no \texttt{submit} do form}

  Troquem o \texttt{axios.post} único por este bloco.

\begin{lstlisting}[basicstyle=\ttfamily\footnotesize]
var dados = {
  nome: document.getElementById("nome").value,
  quantidade: Number(
    document.getElementById("quantidade").value)
};
\end{lstlisting}

  A quantidade vai com \texttt{Number}, porque o input devolve texto.
\end{frame}

\begin{frame}[fragile]{O submit --- POST ou PUT}
  \onde{Ainda no \texttt{submit}}

\begin{lstlisting}[basicstyle=\ttfamily\scriptsize]
if (idEditando === "") {
  axios.post("/api/materiais", dados).then(function () {
    document.getElementById("nome").value = "";
    document.getElementById("quantidade").value = "";
    carregarMateriais();
  });
} else {
  axios.put("/api/materiais/" + idEditando, dados)
    .then(function () {
      idEditando = "";
      document.getElementById("nome").value = "";
      document.getElementById("quantidade").value = "";
      carregarMateriais();
    });
}
\end{lstlisting}

  Depois do PUT, \texttt{idEditando} volta para vazio.\\
  O próximo envio volta a ser cadastro.
\end{frame}

\begin{frame}{Quantidade é número}
  O input devolve texto.\\
  No Java o campo é \texttt{int}.

  \vspace{0.5em}
  Por isso: \texttt{Number(...)}.

  \vspace{0.5em}
  Sem isso, o PUT pode gravar errado
  ou o terminal mostra erro ao converter.
\end{frame}

\begin{frame}{O caminho da edição}
  \begin{enumerate}
    \item A pessoa clica em Editar na Caneta
    \item O form mostra o nome e a quantidade atuais
    \item Ela troca a quantidade e envia
    \item O JS manda \texttt{PUT /api/materiais/1} com JSON
    \item O Controller acha o id 1 e atualiza
    \item O JS chama \texttt{carregarMateriais} de novo
  \end{enumerate}

  \vspace{0.5em}
  No F12, aba Network, deve aparecer um PUT.
\end{frame}

\begin{frame}{Conferindo o PUT}
  \begin{itemize}
    \item [ ] Editar preenche o formulário
    \item [ ] Enviar \textbf{não} cria um material duplicado
    \item [ ] A quantidade nova aparece na lista
    \item [ ] Depois de salvar, o próximo envio volta a cadastrar
    \item [ ] No Network, o pedido é PUT, não POST
  \end{itemize}

  \vspace{0.5em}
  Se nascer um item a mais, o submit ainda está só no POST,
  ou o PUT está dando \texttt{add} na lista.
\end{frame}

% #############################################################################
\section{Buscar um pelo id}

\begin{frame}{Lista inteira e um só material}
  GET \texttt{/api/materiais} devolve \textbf{todos}.

  \vspace{0.5em}
  Às vezes precisamos de \textbf{um}:
  \begin{itemize}
    \item para ver só a Caneta
    \item para conferir a quantidade antes de retirar
  \end{itemize}

  \vspace{0.5em}
  O endereço leva o id:\\
  \texttt{GET /api/materiais/1}
\end{frame}

\begin{frame}[fragile]{GET por id}
  \onde{Projeto \textbf{laboratorio}\\
        \texttt{MaterialController.java}, método novo}

\begin{lstlisting}[language=Java,basicstyle=\ttfamily\footnotesize]
@GetMapping("/api/materiais/{id}")
public Material buscar(@PathVariable int id) {
  for (int i = 0; i < materiais.size(); i++) {
    if (materiais.get(i).getId() == id) {
      return materiais.get(i);
    }
  }
  return null;
}
\end{lstlisting}

  O \texttt{for} é o mesmo do PUT e do DELETE.\\
  Só muda o que fazemos quando achamos: devolver, em vez de
  alterar ou apagar.
\end{frame}

\begin{frame}{Dois GET no mesmo Controller}
  O Spring escolhe pelo endereço:

  \vspace{0.4em}
  \begin{center}
  \begin{tabular}{ll}
    \toprule
    Pedido & Método \\
    \midrule
    \texttt{GET /api/materiais} & \texttt{listar} \\
    \texttt{GET /api/materiais/2} & \texttt{buscar} \\
    \bottomrule
  \end{tabular}
  \end{center}

  \vspace{0.5em}
  O \texttt{\{id\}} é o que diferencia os dois.
\end{frame}

\begin{frame}{Teste no navegador}
  Salvem, reiniciem se precisar, e abram:

  \vspace{0.4em}
  \texttt{http://localhost:8080/api/materiais/2}

  \vspace{0.5em}
  Deve aparecer só o Caderno, em JSON.

  \vspace{0.4em}
  \texttt{/api/materiais} continua mostrando a lista toda.

  \vspace{0.5em}
  Um id que não existe devolve página vazia ou \texttt{null}.
\end{frame}

% #############################################################################
\section{Retirar quantidade}

\begin{frame}{Uma ação pode mudar um campo}
  Até agora o cadastro só cria, lê, troca tudo ou apaga.

  \vspace{0.5em}
  Retirar é outra ação: a pessoa pede uma quantidade,
  e o servidor \textbf{diminui} o que está guardado.

  \vspace{0.5em}
  Exemplo: Caderno tem 5.\\
  Retirar 2 $\rightarrow$ fica 3.
\end{frame}

\begin{frame}{A regra}
  Antes de diminuir, conferimos se dá.

  \vspace{0.5em}
  \begin{enumerate}
    \item Achar o material pelo id
    \item Se a quantidade guardada for \textbf{menor}
          do que a pedida: \textbf{não muda nada}
    \item Se der: subtrai e devolve o material
  \end{enumerate}

  \vspace{0.6em}
  \begin{alertblock}{Não fica negativo}
    Pedir 10 de um item que tem 2 não pode passar de 2
    para $-8$. O servidor recusa.
  \end{alertblock}
\end{frame}

\begin{frame}[fragile]{O pedido de retirar}
  \begin{block}{Pedido}
    \begin{itemize}
      \item Verbo: \texttt{POST}
      \item Endereço: \texttt{/api/materiais/2/retirar}
      \item Corpo:
    \end{itemize}
\begin{lstlisting}[basicstyle=\ttfamily\footnotesize]
{"quantidade": 2}
\end{lstlisting}
  \end{block}

  \vspace{0.3em}
  POST porque é uma ação, não uma troca completa do cadastro.\\
  O id continua na URL.
\end{frame}

\begin{frame}{Quando não pode, o servidor recusa}
  Se a quantidade não alcança, não basta devolver o material
  antigo em silêncio. A tela precisa avisar.

  \vspace{0.5em}
  Vamos lançar um erro HTTP 400:

  \vspace{0.3em}
  \texttt{ResponseStatusException}

  \vspace{0.5em}
  O Axios, nesse caso, cai no \texttt{catch}.\\
  Ali mostramos a mensagem na página.
\end{frame}

\begin{frame}[fragile]{Retirar --- achar e conferir}
  \onde{Projeto \textbf{laboratorio}\\
        \texttt{MaterialController.java}}

  Imports novos:\\
  \texttt{HttpStatus} e \texttt{ResponseStatusException}.

\begin{lstlisting}[language=Java,basicstyle=\ttfamily\footnotesize]
import org.springframework.http.HttpStatus;
import org.springframework.web.server.ResponseStatusException;
\end{lstlisting}
\end{frame}

\begin{frame}[fragile]{MaterialController.java --- retirar}
  \onde{\texttt{MaterialController.java}}

\begin{lstlisting}[language=Java,basicstyle=\ttfamily\scriptsize]
@PostMapping("/api/materiais/{id}/retirar")
public Material retirar(@PathVariable int id,
                        @RequestBody Material pedido) {
  for (int i = 0; i < materiais.size(); i++) {
    Material material = materiais.get(i);
    if (material.getId() == id) {
      if (material.getQuantidade() < pedido.getQuantidade()) {
        throw new ResponseStatusException(
            HttpStatus.BAD_REQUEST, "Quantidade insuficiente");
      }
      material.setQuantidade(
          material.getQuantidade() - pedido.getQuantidade());
      return material;
    }
  }
  return null;
}
\end{lstlisting}
\end{frame}

\begin{frame}{O JSON do retirar}
  O corpo só precisa da quantidade:

  \vspace{0.3em}
  \texttt{\{"quantidade": 2\}}

  \vspace{0.5em}
  O Spring monta um \texttt{Material} e preenche
  só esse campo. O resto fica vazio.

  \vspace{0.5em}
  Nós usamos apenas \texttt{pedido.getQuantidade()}.\\
  Não damos add. Não geramos id. Não trocamos o nome.
\end{frame}

\begin{frame}{O \texttt{if} da regra}
  \texttt{material.getQuantidade() < pedido.getQuantidade()}

  \vspace{0.5em}
  Se for verdade: lança a exceção e \textbf{para}.\\
  A linha do \texttt{setQuantidade} não roda.

  \vspace{0.5em}
  Se for falso: subtrai e devolve o material.
\end{frame}

\begin{frame}[fragile]{Botão Retirar no JS}
  \onde{\texttt{static/js/app.js}, no \texttt{for} de
        \texttt{carregarMateriais}}

\begin{lstlisting}[basicstyle=\ttfamily\scriptsize]
var btnRetirar = document.createElement("button");
btnRetirar.textContent = "Retirar 1";
btnRetirar.setAttribute("data-id", material.id);

btnRetirar.addEventListener("click", function (evento) {
  var id = evento.target.getAttribute("data-id");
  axios.post("/api/materiais/" + id + "/retirar", { quantidade: 1 })
    .then(function () {
      document.getElementById("mensagem").textContent = "";
      carregarMateriais();
    })
    .catch(function () {
      document.getElementById("mensagem").textContent =
        "Quantidade insuficiente";
    });
});
li.appendChild(btnRetirar);
\end{lstlisting}
\end{frame}

\begin{frame}{\texttt{then} e \texttt{catch}}
  \begin{itemize}
    \item \texttt{then} --- o servidor aceitou. Atualizamos a lista.
    \item \texttt{catch} --- o servidor recusou. Mostramos o texto
          no \texttt{<p id="mensagem">}.
  \end{itemize}

  \vspace{0.6em}
  Esse \texttt{<p>} já está no \texttt{index.html}.

  \vspace{0.5em}
  Sem o \texttt{catch}, a recusa aparece só no F12
  e a pessoa não entende o que aconteceu.
\end{frame}

\begin{frame}{O caminho da retirada}
  \begin{enumerate}
    \item Clique em Retirar 1 no Grampeador
    \item JS manda POST com \texttt{\{ quantidade: 1 \}}
    \item Controller acha o id, confere, diminui
    \item A lista recarrega com a quantidade nova
    \item Se já estiver em 0, o próximo clique cai no catch
  \end{enumerate}

  \vspace{0.5em}
  No Network: um POST para \texttt{.../retirar}.\\
  Se recusar, o status é 400.
\end{frame}

\begin{frame}{Conferindo a retirada}
  \begin{itemize}
    \item [ ] Grampeador começa com 2
    \item [ ] Dois cliques em Retirar 1 deixam quantidade 0
    \item [ ] O terceiro clique mostra ``Quantidade insuficiente''
    \item [ ] A quantidade \textbf{não} fica negativa
    \item [ ] Caneta e Caderno não mudam nesse teste
  \end{itemize}
\end{frame}

\begin{frame}{Erros que mais aparecem}
  \begin{itemize}
    \item Porta 8080 ocupada
    \item PUT cria duplicado --- ainda está no POST,
          ou o Controller dá \texttt{add}
    \item Editar preenche sempre o último --- faltou
          \texttt{evento.target}
    \item Quantidade vira texto --- faltou \texttt{Number}
    \item Retirar não avisa na tela --- faltou \texttt{catch}
    \item Import de \texttt{ResponseStatusException} errado
  \end{itemize}
\end{frame}

% #############################################################################
\section{Desafios}

\begin{frame}{Agora é com vocês}
  Continuem no projeto \textbf{laboratorio}.\\
  Mesmos arquivos.

  \vspace{0.5em}
  \begin{enumerate}
    \item Retirar a quantidade digitada
    \item Não cadastrar se o nome estiver vazio
    \item No Editar, buscar o material pelo GET do id
  \end{enumerate}

  \vspace{0.5em}
  Salve $\rightarrow$ reinicie se precisar $\rightarrow$ teste no navegador.\\
  F12, aba Network.
\end{frame}

\begin{frame}{Desafio 1 --- quantidade da retirada}
  \onde{\texttt{index.html} e \texttt{static/js/app.js}}

  Hoje o botão sempre manda 1.

  \vspace{0.4em}
  Coloquem um input na página, por exemplo
  \texttt{id="qtd-retirar"}.

  \vspace{0.4em}
  No clique, leiam esse valor com \texttt{Number}\\
  e mandem no JSON no lugar do 1.

  \vspace{0.4em}
  Testem retirar 3 de um item que tem 2:\\
  tem que recusar.
\end{frame}

\begin{frame}{Desafio 2 --- nome vazio}
  \onde{\texttt{static/js/app.js}, no submit}

  Antes do Axios, se o nome estiver vazio,
  mostrem uma mensagem e deem \texttt{return}.

  \vspace{0.5em}
  O POST nem chega a sair.\\
  A lista permanece igual.
\end{frame}

\begin{frame}[fragile]{Desafio 3 --- Editar busca pelo id}
  \onde{\texttt{static/js/app.js}}

  No clique de Editar, em vez de copiar nome e quantidade
  dos \texttt{data-...}, chamem:

\begin{lstlisting}[basicstyle=\ttfamily\footnotesize]
axios.get("/api/materiais/" + id)
  .then(function (resposta) {
    var material = resposta.data;
    idEditando = material.id;
    document.getElementById("nome").value = material.nome;
    document.getElementById("quantidade").value =
      material.quantidade;
  });
\end{lstlisting}

  Vocês passam a usar o GET por id de verdade.
\end{frame}

\begin{frame}{Extra}
  Se terminarem os três, ainda no \textbf{laboratorio}:

  \vspace{0.4em}
  \begin{itemize}
    \item Não permitir quantidade negativa no cadastro
    \item Mostrar o id na lista
    \item Depois de retirar com sucesso, limpar a mensagem
          de erro antiga
  \end{itemize}
\end{frame}

\begin{frame}{Checklist}
  \begin{itemize}
    \item [ ] Estou no projeto \texttt{laboratorio}?
    \item [ ] PUT não dá \texttt{add} na lista?
    \item [ ] \texttt{idEditando} volta para vazio depois de salvar?
    \item [ ] GET \texttt{/api/materiais/2} devolve um só?
    \item [ ] Retirar recusa quando não há quantidade?
    \item [ ] O \texttt{catch} preenche \texttt{\#mensagem}?
  \end{itemize}
\end{frame}

\begin{frame}{Intervalo}
  \begin{center}
    \Large Retomamos às \textbf{20h20}\\[0.8em]
    \normalsize
    Fechem o \texttt{laboratorio}.\\
    Parem o Spring.\\
    Na segunda parte abrimos o \textbf{StockSales}.
  \end{center}
\end{frame}

% #############################################################################
\section{Parte 2 --- StockSales}

\begin{frame}{Mudamos de projeto}
  \onde{Agora é o projeto \textbf{StockSales}.\\
        O laboratorio já está encerrado.}

  Parem o Spring do laboratorio.\\
  Abram a pasta do StockSales e subam de novo.
\end{frame}

\begin{frame}{Onde o StockSales parou}
  Em produtos, vocês já têm:

  \vspace{0.4em}
  \begin{itemize}
    \item Lista na classe do Controller
    \item GET, POST e DELETE
    \item Formulário na tela e botão Excluir
  \end{itemize}

  \vspace{0.5em}
  Ainda não tem PUT.\\
  Ainda não tem clientes.\\
  Ainda não tem venda.
\end{frame}

\begin{frame}{O que vamos fazer no StockSales}
  \begin{enumerate}
    \item PUT de produto --- o mesmo raciocínio do laboratorio
    \item Cadastro de cliente: GET e POST
    \item Venda de \textbf{um} produto, com baixa de estoque
  \end{enumerate}

  \vspace{0.5em}
  É o que acabamos de treinar, com os nomes da loja:\\
  produto, cliente, estoque.
\end{frame}

\begin{frame}[fragile]{Produto.java --- construtor vazio}
  \onde{Projeto \textbf{StockSales}\\
        \texttt{model/Produto.java}}

  Se ainda não tiver este construtor, acrescentem.
  Mantenham o outro.

\begin{lstlisting}[language=Java,basicstyle=\ttfamily\footnotesize]
public Produto() {
}
\end{lstlisting}

  Sem ele, POST e PUT de produto quebram.
\end{frame}

\begin{frame}[fragile]{PUT de produto}
  \onde{\texttt{controller/ProdutoController.java}}

\begin{lstlisting}[language=Java,basicstyle=\ttfamily\footnotesize]
@PutMapping("/api/produtos/{id}")
public Produto atualizar(@PathVariable int id,
                         @RequestBody Produto dados) {
  for (int i = 0; i < produtos.size(); i++) {
    if (produtos.get(i).getId() == id) {
      produtos.get(i).setNome(dados.getNome());
      produtos.get(i).setPreco(dados.getPreco());
      produtos.get(i).setEstoque(dados.getEstoque());
      return produtos.get(i);
    }
  }
  return null;
}
\end{lstlisting}
\end{frame}

\begin{frame}{Tela de produtos --- Editar}
  \onde{\texttt{static/js/produtos.js}}

  Copiem o esquema do laboratorio:

  \vspace{0.4em}
  \begin{itemize}
    \item Variável \texttt{idEditando}
    \item Botão Editar com \texttt{data-id}
    \item Submit: se o id estiver vazio $\rightarrow$ POST;
          se tiver id $\rightarrow$ PUT
    \item Preço e estoque com \texttt{Number}
  \end{itemize}

  \vspace{0.4em}
  No HTML, o form de produtos que vocês já usam no POST
  serve para o PUT.
\end{frame}

\begin{frame}{Três cadastros, uma loja}
  Daqui a pouco teremos:

  \vspace{0.4em}
  \begin{itemize}
    \item produtos
    \item clientes
    \item vendas
  \end{itemize}

  \vspace{0.5em}
  Cada um pode ter o seu Controller.\\
  As listas precisam ser as \textbf{mesmas}, senão a venda
  não acha o produto e o estoque não baixa.

  \vspace{0.5em}
  Por isso as listas saem dos Controllers
  e vão para uma classe só: \texttt{Dados}.
\end{frame}

\begin{frame}[fragile]{Classe Dados}
  \onde{Projeto \textbf{StockSales}\\
        Arquivo novo:\\
        \texttt{src/main/java/com/entra21/stocksales/Dados.java}}

\begin{lstlisting}[language=Java,basicstyle=\ttfamily\footnotesize]
package com.entra21.stocksales;

import java.util.ArrayList;
import com.entra21.stocksales.model.Produto;

public class Dados {
  public static ArrayList<Produto> produtos = new ArrayList<>();
  public static int proximoIdProduto = 1;
}
\end{lstlisting}

  \texttt{static} = uma lista só, visível para qualquer classe
  do projeto.
\end{frame}

\begin{frame}[fragile]{Mover a lista de produtos}
  \onde{\texttt{ProdutoController.java}}

  Tirem os campos \texttt{produtos} e \texttt{proximoId}
  da classe.

  \vspace{0.3em}
  O construtor passa a gravar em \texttt{Dados}:

\begin{lstlisting}[language=Java,basicstyle=\ttfamily\scriptsize]
public ProdutoController() {
  if (Dados.produtos.isEmpty()) {
    Dados.produtos.add(new Produto(1, "Teclado Mecanico", 250, 12));
    Dados.produtos.add(new Produto(2, "Mouse Gamer", 120, 30));
    Dados.produtos.add(new Produto(3, "Monitor 24", 900, 8));
    Dados.produtos.add(new Produto(4, "Headset", 200, 15));
    Dados.proximoIdProduto = 5;
  }
}
\end{lstlisting}
\end{frame}

\begin{frame}[fragile]{GET, POST, PUT e DELETE usam Dados}
  \onde{\texttt{ProdutoController.java}}

  Troquem \texttt{produtos} por \texttt{Dados.produtos}.\\
  Troquem \texttt{proximoId} por \texttt{Dados.proximoIdProduto}.

\begin{lstlisting}[language=Java,basicstyle=\ttfamily\footnotesize]
@GetMapping("/api/produtos")
public ArrayList<Produto> listarProdutos() {
  return Dados.produtos;
}
\end{lstlisting}

  \vspace{0.4em}
  O \texttt{if (Dados.produtos.isEmpty())} evita duplicar
  o seed se o Spring criar o Controller de novo.
\end{frame}

\begin{frame}[fragile]{Classe Cliente}
  \onde{Arquivo novo:\\
        \texttt{src/main/java/com/entra21/stocksales/model/Cliente.java}}

\begin{lstlisting}[language=Java,basicstyle=\ttfamily\scriptsize]
public class Cliente {
  private int id;
  private String nome;
  private String email;
  private String telefone;

  public Cliente() {}

  public Cliente(int id, String nome, String email, String telefone) {
    this.id = id;
    this.nome = nome;
    this.email = email;
    this.telefone = telefone;
  }
}
\end{lstlisting}

  Completem os gets e sets dos quatro atributos.
\end{frame}

\begin{frame}[fragile]{Listas de cliente em Dados}
  \onde{\texttt{Dados.java}, acrescentar}

\begin{lstlisting}[language=Java,basicstyle=\ttfamily\footnotesize]
public static ArrayList<Cliente> clientes = new ArrayList<>();
public static int proximoIdCliente = 1;
\end{lstlisting}

  Import de \texttt{Cliente} no topo do arquivo.
\end{frame}

\begin{frame}[fragile]{ClienteController --- seed e GET}
  \onde{Arquivo novo:\\
        \texttt{controller/ClienteController.java}}

\begin{lstlisting}[language=Java,basicstyle=\ttfamily\scriptsize]
@RestController
public class ClienteController {

  public ClienteController() {
    if (Dados.clientes.isEmpty()) {
      Dados.clientes.add(new Cliente(1, "Ana", "ana@email.com", "1199"));
      Dados.clientes.add(new Cliente(2, "Bruno", "bruno@email.com", "1188"));
      Dados.clientes.add(new Cliente(3, "Carla", "carla@email.com", "1177"));
      Dados.proximoIdCliente = 4;
    }
  }

  @GetMapping("/api/clientes")
  public ArrayList<Cliente> listar() {
    return Dados.clientes;
  }
}
\end{lstlisting}
\end{frame}

\begin{frame}[fragile]{POST de cliente}
  \onde{Ainda em \texttt{ClienteController.java}}

\begin{lstlisting}[language=Java,basicstyle=\ttfamily\footnotesize]
@PostMapping("/api/clientes")
public Cliente cadastrar(@RequestBody Cliente cliente) {
  cliente.setId(Dados.proximoIdCliente);
  Dados.proximoIdCliente = Dados.proximoIdCliente + 1;
  Dados.clientes.add(cliente);
  return cliente;
}
\end{lstlisting}

  Testem \texttt{/api/clientes} no navegador:\\
  Ana, Bruno e Carla.
\end{frame}

\begin{frame}[fragile]{Tela de clientes}
  \onde{Arquivos novos:\\
        \texttt{static/clientes.html}\\
        \texttt{static/js/clientes.js}}

  Copiem o desenho de produtos:

  \vspace{0.3em}
  \begin{itemize}
    \item Formulário: nome, e-mail, telefone
    \item Lista na \texttt{<ul>}
    \item \texttt{carregarClientes} com GET
    \item Submit com POST e \texttt{preventDefault}
  \end{itemize}

  \vspace{0.4em}
  Na home, um link para \texttt{/clientes.html}.
\end{frame}

\begin{frame}{O que é uma venda nesta aula}
  Uma venda, hoje, tem \textbf{um} cliente e \textbf{um} produto.

  \vspace{0.5em}
  \begin{itemize}
    \item Quem compra: o cliente
    \item O quê: um produto
    \item Quantas unidades
    \item Total = preço $\times$ quantidade
  \end{itemize}

  \vspace{0.5em}
  Ao confirmar: o estoque daquele produto diminui.\\
  Se não houver estoque, a venda \textbf{não} é gravada.
\end{frame}

\begin{frame}[fragile]{Classe Venda}
  \onde{Arquivo novo:\\
        \texttt{src/main/java/com/entra21/stocksales/model/Venda.java}}

\begin{lstlisting}[language=Java,basicstyle=\ttfamily\footnotesize]
public class Venda {
  private int id;
  private int clienteId;
  private int produtoId;
  private int quantidade;
  private double total;

  public Venda() {
  }
}
\end{lstlisting}

  Construtor vazio + gets e sets dos cinco atributos.
\end{frame}

\begin{frame}[fragile]{Vendas em Dados}
  \onde{\texttt{Dados.java}}

\begin{lstlisting}[language=Java,basicstyle=\ttfamily\footnotesize]
public static ArrayList<Venda> vendas = new ArrayList<>();
public static int proximoIdVenda = 1;
\end{lstlisting}
\end{frame}

\begin{frame}[fragile]{O POST da venda}
  O JavaScript manda só isto:

\begin{lstlisting}[basicstyle=\ttfamily\footnotesize]
{"clienteId": 1, "produtoId": 2, "quantidade": 3}
\end{lstlisting}

  \vspace{0.4em}
  O servidor faz o resto:
  \begin{enumerate}
    \item Achar o cliente
    \item Achar o produto
    \item Se \texttt{estoque < quantidade}, recusar
    \item \texttt{total = preco * quantidade}
    \item Diminuir o estoque
    \item Guardar a venda e devolver o JSON
  \end{enumerate}
\end{frame}

\begin{frame}[fragile]{Começar o VendaController}
  \onde{Arquivo novo:\\
        \texttt{controller/VendaController.java}}

\begin{lstlisting}[language=Java,basicstyle=\ttfamily\footnotesize]
@RestController
public class VendaController {
\end{lstlisting}

  Os métodos GET e POST entram \textbf{dentro} desta classe.\\
  O POST é longo: vamos escrever em três passos.
\end{frame}

\begin{frame}[fragile]{VendaController --- achar cliente}
  \onde{\texttt{VendaController.java}, método \texttt{cadastrar}}

\begin{lstlisting}[language=Java,basicstyle=\ttfamily\footnotesize]
@PostMapping("/api/vendas")
public Venda cadastrar(@RequestBody Venda pedido) {
  Cliente cliente = null;
  for (int i = 0; i < Dados.clientes.size(); i++) {
    if (Dados.clientes.get(i).getId() == pedido.getClienteId()) {
      cliente = Dados.clientes.get(i);
    }
  }
  if (cliente == null) {
    throw new ResponseStatusException(
        HttpStatus.BAD_REQUEST, "Cliente nao encontrado");
  }
\end{lstlisting}

  O método ainda não acabou. Continua no próximo slide.
\end{frame}

\begin{frame}[fragile]{VendaController --- achar produto}
  \onde{Ainda no mesmo método \texttt{cadastrar}}

\begin{lstlisting}[language=Java,basicstyle=\ttfamily\footnotesize]
  Produto produto = null;
  for (int i = 0; i < Dados.produtos.size(); i++) {
    if (Dados.produtos.get(i).getId() == pedido.getProdutoId()) {
      produto = Dados.produtos.get(i);
    }
  }
  if (produto == null) {
    throw new ResponseStatusException(
        HttpStatus.BAD_REQUEST, "Produto nao encontrado");
  }
\end{lstlisting}
\end{frame}

\begin{frame}[fragile]{VendaController --- estoque e gravação}
  \onde{Ainda no mesmo método \texttt{cadastrar}}

\begin{lstlisting}[language=Java,basicstyle=\ttfamily\footnotesize]
  if (produto.getEstoque() < pedido.getQuantidade()) {
    throw new ResponseStatusException(
        HttpStatus.BAD_REQUEST, "Estoque insuficiente");
  }

  pedido.setId(Dados.proximoIdVenda);
  Dados.proximoIdVenda = Dados.proximoIdVenda + 1;
  pedido.setTotal(produto.getPreco() * pedido.getQuantidade());
  produto.setEstoque(
      produto.getEstoque() - pedido.getQuantidade());
  Dados.vendas.add(pedido);
  return pedido;
}
\end{lstlisting}
\end{frame}

\begin{frame}[fragile]{GET das vendas}
  \onde{Ainda em \texttt{VendaController.java}}

\begin{lstlisting}[language=Java,basicstyle=\ttfamily\footnotesize]
@GetMapping("/api/vendas")
public ArrayList<Venda> listar() {
  return Dados.vendas;
}
\end{lstlisting}

  \vspace{0.4em}
  Não esqueçam os imports:
  \texttt{HttpStatus}, \texttt{ResponseStatusException},
  \texttt{Dados}, \texttt{Cliente}, \texttt{Produto}, \texttt{Venda}.
\end{frame}

\begin{frame}[fragile]{A tela precisa de um \texttt{select}}
  Em vez de a pessoa lembrar o id, a página mostra uma lista
  para escolher.

  \vspace{0.4em}
\begin{lstlisting}[basicstyle=\ttfamily\footnotesize]
<select id="clienteId"></select>
\end{lstlisting}

  \vspace{0.4em}
  Cada \texttt{<option>} tem o \texttt{value} = id
  e o texto = nome.

  \vspace{0.5em}
  O JS preenche as opções com o GET de clientes
  e o GET de produtos.
\end{frame}

\begin{frame}[fragile]{vendas-nova.html}
  \onde{\texttt{static/vendas-nova.html}, arquivo novo}

\begin{lstlisting}[basicstyle=\ttfamily\scriptsize]
<h1>Nova venda</h1>
<a href="/">Voltar</a>
<label>Cliente</label>
<select id="clienteId"></select>
<label>Produto</label>
<select id="produtoId"></select>
<label>Quantidade</label>
<input id="quantidade" type="number">
<button id="btn-vender" type="button">Confirmar venda</button>
<p id="mensagem"></p>
<script src="https://cdn.jsdelivr.net/npm/axios/dist/axios.min.js"></script>
<script src="/js/vendas-nova.js"></script>
\end{lstlisting}
\end{frame}

\begin{frame}[fragile]{Preencher os selects}
  \onde{\texttt{static/js/vendas-nova.js}, arquivo novo}

\begin{lstlisting}[basicstyle=\ttfamily\footnotesize]
function carregarClientes() {
  axios.get("/api/clientes").then(function (resposta) {
    var select = document.getElementById("clienteId");
    select.innerHTML = "";
    var clientes = resposta.data;
    for (var i = 0; i < clientes.length; i++) {
      var opcao = document.createElement("option");
      opcao.value = clientes[i].id;
      opcao.textContent = clientes[i].nome;
      select.appendChild(opcao);
    }
  });
}
\end{lstlisting}

  O de produtos é igual, com nome e estoque no texto.
\end{frame}

\begin{frame}[fragile]{Select de produtos}
  \onde{Ainda em \texttt{vendas-nova.js}}

\begin{lstlisting}[basicstyle=\ttfamily\footnotesize]
function carregarProdutos() {
  axios.get("/api/produtos").then(function (resposta) {
    var select = document.getElementById("produtoId");
    select.innerHTML = "";
    var produtos = resposta.data;
    for (var i = 0; i < produtos.length; i++) {
      var opcao = document.createElement("option");
      opcao.value = produtos[i].id;
      opcao.textContent = produtos[i].nome
        + " --- estoque: " + produtos[i].estoque;
      select.appendChild(opcao);
    }
  });
}
\end{lstlisting}
\end{frame}

\begin{frame}[fragile]{Confirmar a venda --- o JSON}
  \onde{Ainda em \texttt{vendas-nova.js}}

\begin{lstlisting}[basicstyle=\ttfamily\footnotesize]
document.getElementById("btn-vender")
  .addEventListener("click", function () {
    var dados = {
      clienteId: Number(document.getElementById("clienteId").value),
      produtoId: Number(document.getElementById("produtoId").value),
      quantidade: Number(document.getElementById("quantidade").value)
    };
\end{lstlisting}

  Os três valores vão como número, não como texto.
\end{frame}

\begin{frame}[fragile]{Confirmar a venda --- enviar}
  \onde{Ainda no clique de Confirmar}

\begin{lstlisting}[basicstyle=\ttfamily\footnotesize]
    axios.post("/api/vendas", dados)
      .then(function () {
        document.getElementById("mensagem").textContent =
          "Venda registrada";
        carregarProdutos();
      })
      .catch(function () {
        document.getElementById("mensagem").textContent =
          "Nao foi possivel registrar a venda";
      });
  });

carregarClientes();
carregarProdutos();
\end{lstlisting}
\end{frame}

\begin{frame}[fragile]{Histórico de vendas}
  \onde{\texttt{static/vendas.html} e \texttt{static/js/vendas.js}}

  Página com uma \texttt{<ul id="lista-vendas">}.

\begin{lstlisting}[basicstyle=\ttfamily\scriptsize]
axios.get("/api/vendas").then(function (resposta) {
  var vendas = resposta.data;
  var lista = document.getElementById("lista-vendas");
  lista.innerHTML = "";
  for (var i = 0; i < vendas.length; i++) {
    var li = document.createElement("li");
    li.textContent = "Cliente " + vendas[i].clienteId
      + " - produto " + vendas[i].produtoId
      + " - qtd " + vendas[i].quantidade
      + " - total R$ " + vendas[i].total;
    lista.appendChild(li);
  }
});
\end{lstlisting}

  Na home: links para nova venda e para o histórico.
\end{frame}

\begin{frame}{O caminho da venda}
  \begin{enumerate}
    \item A pessoa escolhe Ana, Mouse Gamer, quantidade 2
    \item JS manda POST \texttt{/api/vendas}
    \item Servidor acha cliente e produto nas listas de
          \texttt{Dados}
    \item Confere estoque, calcula total, diminui estoque
    \item A tela de produtos, ao recarregar, já mostra
          o estoque menor
  \end{enumerate}
\end{frame}

\begin{frame}{Conferindo no StockSales}
  \begin{itemize}
    \item [ ] Editar produto não duplica
    \item [ ] \texttt{/api/clientes} lista Ana, Bruno e Carla
    \item [ ] Dá para cadastrar um cliente novo
    \item [ ] Nova venda baixa o estoque
    \item [ ] Vender acima do estoque mostra a mensagem
    \item [ ] Histórico lista a venda com o total
  \end{itemize}
\end{frame}

\begin{frame}{Extra no StockSales}
  Se der tempo:

  \vspace{0.4em}
  \begin{itemize}
    \item PUT e DELETE de cliente
    \item No histórico, mostrar o \textbf{nome} do cliente
          e do produto, não só o id. Percorram as listas
          de \texttt{Dados} no JS, ou incluam os nomes
          na resposta da venda
    \item Não permitir quantidade 0 ou negativa
  \end{itemize}
\end{frame}

\begin{frame}{O que não entra hoje}
  \begin{itemize}
    \item Vários produtos na mesma venda
    \item Cancelar venda e devolver estoque
    \item Banco de dados
    \item Login
  \end{itemize}

  \vspace{0.5em}
  \begin{block}{Meta no StockSales}
    PUT de produto.\\
    Clientes cadastráveis.\\
    Uma venda de um item, com estoque baixando.
  \end{block}
\end{frame}

\begin{frame}{Critérios de aceite}
  \begin{itemize}
    \item [ ] Produto pode ser alterado
    \item [ ] Cliente aparece na lista e no select da venda
    \item [ ] Venda grava total no servidor
    \item [ ] Estoque diminui
    \item [ ] Estoque insuficiente recusa a venda inteira
    \item [ ] Home leva a produtos, clientes, nova venda
              e histórico
  \end{itemize}
\end{frame}

% #############################################################################
\section{Recapitulando}

\begin{frame}{O que vimos hoje}
  \begin{enumerate}
    \item PUT altera um item que já existe, pelo id
    \item GET com id devolve um só
    \item Retirar / vender é uma ação com regra
    \item Se a regra falhar, o servidor recusa e o JS usa catch
    \item No StockSales, as listas ficam em \texttt{Dados}
          para todo mundo ver a mesma loja
  \end{enumerate}
\end{frame}

\begin{frame}{Na próxima aula}
  \begin{block}{Aula 05}
    \begin{itemize}
      \item Organizar o código: as listas saem do
            \texttt{static} de \texttt{Dados}
      \item Validação
      \item Erros HTTP com mais clareza
    \end{itemize}
  \end{block}

  \vspace{0.6em}
  Tragam o StockSales com venda de um item funcionando.
\end{frame}

\begin{frame}
  \begin{center}
    \Huge Obrigado!\\[1.2em]
    \normalsize Mauricio Duailibi Neto\\
    Turma Java Entra21
  \end{center}
\end{frame}

\end{document}
